\documentclass[12pt]{article}
\usepackage[utf8]{inputenc}
\usepackage[T1]{fontenc}
\usepackage{hyperref}
\title{Yazghulami Language Resources for K-12 Curriculum Development}
\author{Compiled by Claude Sonnet 5.5 (claude-sonnet-5-5) for Andr\'as Kornai}
\date{}
\begin{document}
\maketitle

\section{Introduction}
Yazghulami (Yazgulyam; ISO 639-3 yah) is a southeastern Iranian language of the Pamir group, spoken along the Yazghulom River in Gorno-Badakhshan, Tajikistan. Speaker counts vary: about 9,000 in Wikipedia, about 4,000 (1994) in an older Rosetta entry. Speakers are bilingual in Tajik. Wikipedia reports that the language is not traditionally written, but that linguists in Russia and Tajikistan recently developed an alphabet. No school materials were found.

\section{Primary Educational Resources}
No primers, textbooks or school programmes were found in the search results.

\section{Linguistic and Reference Materials}
\begin{itemize}
\item Wikipedia, Yazghulami language. \url{https://en.wikipedia.org/wiki/Yazghulami_language}
\item Webonary Yazghulami dictionary, with grammar, phonology and bibliography sections in its menu. The contents could not be checked from the snippet. \url{https://www.webonary.work/yazghulami/?p=146825}
\item Lemontree typological database entry for Yazghulami, drawing on the Russian encyclopedia of Iranian languages. \url{https://lemontree210.pythonanywhere.com/en/doculect/yazghulami}
\item Glosbe Yazghulami--English page, which looked sparse. \url{https://glosbe.com/yah/en}
\end{itemize}

\section{Religious and Community Resources}
None were found in the search results.

\section{Digital and Audio Resources}
Catalogue pages: Endangered Languages Project, \url{https://endangeredlanguages.com/lang/yah?hl=pt-br}, bibliography at \url{https://endangeredlanguages.com/node/38718/bibliography}; Rosetta Project, \url{https://rosettapanglossia.longnow.org/index.php?title=Yah}. Their audio content was not checked.

\section{Curriculum Development Context}
Little ready-to-use educational material was found. A curriculum would start from the new alphabet, which needs confirming, and from the Webonary dictionary and grammar. Tajik and Russian sources would need further searching.
\end{document}